\documentclass[11pt]{article}
\input{../../shared/project_style.tex}
\rhead{Basics 7 Textbook Draft}

\begin{document}
\DocTitle{Basics 7: Fourier Analysis and Normal Modes}{I - Mathematical and computational foundations}{Spatial harmonics, wave numbers, dispersion relations, toroidal mode labels, and coupling}

\begin{overviewbox}
\textbf{Textbook draft, version 1.0.} Fourier analysis expresses fields as sums of simple oscillations. In toroidal plasmas, poloidal and toroidal harmonics reveal how geometry couples wave components. This chapter is designed for a reader with no prior plasma-physics training. It distinguishes exact identities from physical approximations and connects the central equations to later simulation modules.
\end{overviewbox}

\ProjectTOC

\section{Learning objectives}
After completing this note, the reader should be able to:
\begin{itemize}[leftmargin=1.6em]
\item Derive Fourier-series coefficients from orthogonality.
\item Interpret Fourier transforms and wave number.
\item Use complex exponentials to represent real fields.
\item Derive normal-mode dispersion relations.
\item Interpret poloidal and toroidal mode numbers.
\item Recognize aliasing, spectral resolution, and harmonic coupling.
\end{itemize}

\section{Prerequisites and reading strategy}
\begin{itemize}[leftmargin=1.6em]
\item Basics 2, 5, and 6.
\end{itemize}
On a first reading, emphasize physical meaning and dimensional checks. On a second reading, reproduce the derivations. Attempt the computational laboratory only after the limiting cases are understood.

\section{Fourier series and orthogonality}

For a $2\pi$-periodic field,
\begin{equation}
f(\theta)=\sum_{m=-\infty}^{\infty}\hat f_me^{im\theta}.
\end{equation}
Orthogonality gives
\begin{equation}
\hat f_m=\frac{1}{2\pi}\int_0^{2\pi}f(\theta)e^{-im\theta}\,d\theta.
\end{equation}
A real field obeys $\hat f_{-m}=\hat f_m^*$. Parseval relation connects physical-space and spectral energy:
\begin{equation}
\frac{1}{2\pi}\int_0^{2\pi}|f|^2d\theta=\sum_m|\hat f_m|^2.
\end{equation}


\section{Fourier transforms and differentiation}

For a nonperiodic coordinate,
\begin{align}
\hat f(k)&=\int_{-\infty}^{\infty}f(x)e^{-ikx}\,dx,\\
f(x)&=\frac{1}{2\pi}\int_{-\infty}^{\infty}\hat f(k)e^{ikx}\,dk.
\end{align}
Derivatives become multiplication:
\begin{equation}
\widehat{f_x}=ik\hat f,\qquad
\widehat{f_{xx}}=-k^2\hat f.
\end{equation}
This converts constant-coefficient linear PDEs into algebraic equations for each wave number.


\section{Normal modes and dispersion}

Assume
\begin{equation}
u(\mathbf x,t)=\tilde u\,e^{i(\mathbf k\cdot\mathbf x-\omega t)}.
\end{equation}
Then $\partial_t\to-i\omega$ and $\nabla\to i\mathbf k$. A nonzero amplitude requires a dispersion relation
\begin{equation}
D(\omega,\mathbf k)=0.
\end{equation}
For the wave equation, $\omega^2=c^2k^2$. Phase velocity is $\omega/k$ and group velocity is $d\omega/dk$.


\section{Poloidal and toroidal harmonics}

On a torus,
\begin{equation}
f(s,\theta,\phi,t)=\sum_{m,n}\hat f_{mn}(s)
e^{i(m\theta-n\phi-\omega t)}.
\end{equation}
$m$ is the poloidal mode number and $n$ toroidal under this convention. Other sign conventions exist. The radial amplitude is not generally sinusoidal because equilibrium quantities and boundaries vary with $s$.


\section{Geometric harmonic coupling}

In a uniform medium, Fourier harmonics can evolve independently. A coefficient
\begin{equation}
C(\theta)=C_0+\epsilon\cos\theta
\end{equation}
multiplied by $e^{im\theta}$ produces $m-1$, $m$, and $m+1$ components. Toroidicity therefore couples neighboring poloidal harmonics. Three-dimensional stellarator geometry contains many equilibrium harmonics and generates broader coupling networks.


\section{Sampling, aliasing, and spectral interpretation}

With $N$ samples, only modes below the Nyquist limit are represented uniquely. Higher modes appear as aliases. Nonlinear products generate harmonics above the original bandwidth, motivating dealiasing in spectral solvers.

A physical eigenmode can contain many Fourier harmonics. The labels $(m,n)$ describe components, not necessarily distinct independent modes.

\begin{physicsbox}
Always record the sign convention and array ordering for Fourier coefficients. A silent convention mismatch can reverse $k_\parallel$ and corrupt resonance analysis.
\end{physicsbox}


\section{Computational laboratory}

Sample a signal containing $m=3$ and $m=7$, compute its discrete Fourier transform, and recover amplitudes and phases. Reduce the grid to demonstrate aliasing. Multiply a pure harmonic by $1+\epsilon\cos\theta$ and verify generation of neighboring harmonics. Compare Python and MATLAB FFT ordering.


\section{Summary}
\begin{itemize}[leftmargin=1.6em]
\item Fourier analysis decomposes periodic fields into orthogonal harmonics.
\item Normal-mode substitution reduces linear PDEs to dispersion or eigenvalue problems.
\item Toroidal fields use poloidal and toroidal mode numbers.
\item Geometry couples harmonics through nonuniform coefficients.
\item Sampling and conventions determine whether spectra are interpreted correctly.
\end{itemize}

\SymbolsAcronymsSection
\begin{symtable}
$m,n$ & mode numbers & Poloidal and toroidal indices. \\
$\hat f_m$ & Fourier coefficient & Complex amplitude of harmonic $m$. \\
$\mathbf k$ & wave vector & Spatial phase gradient. \\
$\omega$ & angular frequency & Temporal phase rate. \\
$v_p,v_g$ & velocities & Phase and group velocity. \\
\end{symtable}

\section{Exercises}
\begin{enumerate}[leftmargin=1.8em]
\item Derive the Fourier coefficient formula from orthogonality.
\item Prove the derivative property of the Fourier transform.
\item Derive phase and group velocity for $\omega^2=\omega_0^2+c^2k^2$.
\item Interpret constant-phase surfaces for an $(m,n)$ toroidal harmonic.
\item Explain the conjugate symmetry of a real field.
\end{enumerate}

\section{References}
\begin{thebibliography}{99}
\bibitem{ref1} G. B. Arfken, H. J. Weber, and F. E. Harris, \emph{Mathematical Methods for Physicists}, 7th ed., Academic Press (2013).
\bibitem{ref2} G. Strang, \emph{Linear Algebra and Its Applications}, 4th ed., Brooks/Cole (2006).
\bibitem{ref3} W. A. Strauss, \emph{Partial Differential Equations: An Introduction}, 2nd ed., Wiley (2007).
\bibitem{ref4} T. H. Stix, \emph{Waves in Plasmas}, American Institute of Physics (1992).
\bibitem{ref5} D. A. Gurnett and A. Bhattacharjee, \emph{Introduction to Plasma Physics}, 2nd ed., Cambridge University Press (2017).
\bibitem{ref6} J. P. Freidberg, \emph{Ideal MHD}, Cambridge University Press (2014).
\bibitem{ref7} J. P. Goedbloed, R. Keppens, and S. Poedts, \emph{Advanced Magnetohydrodynamics}, Cambridge University Press (2010).
\bibitem{ref8} H. Goedbloed and S. Poedts, \emph{Principles of Magnetohydrodynamics}, Cambridge University Press (2004).
\end{thebibliography}

\end{document}
